\section{Acquisitions with persistent hidden state}\label{sec:v5-dependent}
The preceding realization can be used when acquisition does not
regenerate the physical state. The appropriate comparison concerns
raw prediction risk: the full-history posterior need not equal the
posterior based on the latest observation alone.

Let $E$ and $\mathsf Y$ be standard Borel spaces, let $T:E\to E$
be Borel, let $f:E\to\R^r$ be Borel with $\|f\|\le F$, and let
$\mu$ be a probability with $T_\#\mu=\mu$. Acquisitions occur at
$\tau_0=0$ and $\tau_{j+1}=\tau_j+L_j$, where the $L_j$ are
independent positive integer random variables of common finite
mean $\ell$. The acquisition clock is independent of the hidden
transition and observation randomizations. Between acquisitions
$X_t=TX_{t-1}$. At an acquisition $t>0$ the transition and
observation are specified, in this order, by
\begin{equation}\label{eq:v5-dependentkernel}
 X_t\mid X_{t-1}=x\sim Q(x,du),\qquad
 Y_t\mid X_t=u\sim K(u,dy).
\end{equation}
At time zero $X_0\sim\mu$ and $Y_0\sim K(X_0,\cdot)$.
We assume
\begin{equation}\label{eq:v5-minorization}
 \mu Q=\mu,\qquad Q(x,A)\ge\epsilon\mu(A)
 \quad(x\in E,\ A\text{ Borel}),\qquad 0<\epsilon\le1.
\end{equation}
The whole state $X_{t-1}$ can influence $X_t$ through the remaining
part of $Q$. No regeneration indicator is supplied to the observer.
The observer sees the acquisition indicators and the $Y_t$ only
when acquisitions occur; the persistent-state convention is that
of Section~\ref{sec:v5-model}.

Use a single pair $(U,Y)\sim\mu(du)K(u,dy)$ to define
\begin{align}
 s_k&=\Pp(L_0>k),\quad w_k=s_k/\ell,\quad
 W(n)=\sum_{k\ge n}w_k,\nonumber\\
 m_k(y)&=\E[f(T^kU)\mid Y=y],\quad
 Z(y)=(\sqrt{w_k}m_k(y))_{k\ge0},\label{eq:v5-dependentorbit}\\
 B_0&=\sum_{k\ge0}w_k\E\|f(T^kU)-m_k(Y)\|^2.\nonumber
\end{align}
Let $R_M^{\mathrm{dep}}$ denote the infimum of limsup average
\emph{raw} squared loss in this dependent experiment.

\begin{theorem}[Minorized acquisition comparison]\label{thm:v5-dependent}
Under \eqref{eq:v5-dependentkernel}--\eqref{eq:v5-minorization},
\begin{equation}\label{eq:v5-dependentlower}
 R_M^{\mathrm{dep}}\ge\epsilon\bigl(B_0+e_M^2(Z)\bigr)
 \quad(M\ge1).
\end{equation}
The same lower bound holds for the liminf average loss of each
randomized, time-dependent $M$-state observer. For every $N,n\ge1$
with $Nn+1\le M$,
\begin{equation}\label{eq:v5-dependentupper}
 R_M^{\mathrm{dep}}\le B_0+e_N^2(Z^{<n})+F^2W(n).
\end{equation}
A one-state zero decoder always has risk at most $F^2$.
More generally, every observer that overwrites its entire state
using only the latest observation at each acquisition has exactly
the same sequence of expected one-time losses as in the
independent-acquisition experiment with law $\mu K$ and the same
clock, provided it is initialized in the same way.
\end{theorem}

\begin{proof}
Condition on the entire acquisition-indicator sequence. At time
zero the hidden marginal is $\mu$. Each intervening transition is
either $T$ or $Q$, both of which preserve $\mu$. Induction therefore
shows that the hidden state at each acquisition has marginal
$\mu$, even after this conditioning. Its observation pair has
law $\mu K$. Future survival for $k$ steps has probability $s_k$,
independent of that pair. Thus an observer which restarts from
the current observation has the same age-dependent expected loss
as in the independent experiment, and hence the same expected
one-time loss after averaging over the common clock. This proves
the last assertion. Applying the $Nn+1$-state construction of
Theorem~\ref{thm:v4-orbit} proves
\eqref{eq:v5-dependentupper}; the constant decoder proves the
remaining upper bound.

For the converse fix $n$. At an acquisition condition on the
entire hidden and observed past immediately before the transition,
and use the independent uniform representation of the observer's
kernels from Lemma~\ref{lem:v4-kernels}. Freeze its coding uniforms,
including those used on the hypothetical next $n$ nonacquisition
steps. The $M$ possible post-update states determine at most $M$
length-$n$ decoder strings. Conditional on this past, the new
observation pair has law $Q(x,du)K(u,dy)$, which dominates
$\epsilon\mu(du)K(u,dy)$. Squared loss is nonnegative, so its
expected contribution up to the next acquisition, truncated at
$n$, is at least
\[
 \epsilon\inf_{i(\cdot),\,c_{i,k}}
 \sum_{k<n}s_k\E_{\mu K}\|f(T^kU)-c_{i(Y),k}\|^2
 =\epsilon\ell\bigl(B_{0,<n}+e_M^2(Z^{<n})\bigr).
\]
For the equality, condition on $Y$ and use orthogonality around
$m_k(Y)$; the remaining infimum is precisely finite-dimensional
quantization. Enlarging the possible codebooks to all vectors
only weakens this lower bound. The encoder is permitted to depend
on the conditioned past, but the resulting bound does not.

For a horizon $H$, sum these truncated contributions only over
acquisitions at times at most $H-n$. Each contribution stops at
the next acquisition, so distinct terms charge disjoint time
indices, path by path. The expected number of these acquisitions,
divided by $H$, converges to $1/\ell$ by the elementary renewal
theorem. Taking liminf in $H$ proves the lower bound with
$B_{0,<n}$ and $Z^{<n}$. Finally let $n\to\infty$.
The noise terms converge by boundedness. The quantization terms
increase to $e_M^2(Z)$: projection gives one inequality, and
extending a truncated codebook by zero gives
\[
 e_M^2(Z)\le e_M^2(Z^{<n})+\E\|Z-Z^{<n}\|^2
 \le e_M^2(Z^{<n})+F^2W(n).
\]
This proves \eqref{eq:v5-dependentlower}.
\end{proof}

The reference number $B_0$ in this theorem is \emph{not} asserted
to be the Bayes risk conditional on the entire dependent history.
Older observations can improve that posterior. The theorem uses
minorization to obtain a lower bound and invariant marginals to
obtain an upper bound; it does not replace a dependent conditional
law by an independent one.

\begin{corollary}[The pressure law survives partial refreshment]
\label{cor:v5-dependentpressure}
Take the repeller, Gibbs law and geometric acquisition clock of
Section~\ref{sec:v5-model}, but use any kernel $Q$ satisfying
\eqref{eq:v5-minorization} at acquisitions. Observe the current
state exactly at those times. Then, for each fixed $\epsilon>0$,
\[
 R_M^{\mathrm{dep}}\asymp G_M(q),\qquad
 \lim_{M\to\infty}\frac{-\log R_M^{\mathrm{dep}}}{\log M}
 =\frac{1-s_*(q)}{s_*(q)}.
\]
The same order holds for supremum one-time loss. An explicit
example is
\begin{equation}\label{eq:v5-partialreset}
 Q(x,du)=(1-\epsilon)\delta_{Tx}(du)+\epsilon\mu(du).
\end{equation}
For $\epsilon<1$ the physical state is not regenerated at every
acquisition. The phase boundary is governed by the acquisition
survival probability $q$, not by the unobserved physical-refresh
survival probability $1-\epsilon(1-q)$.
\end{corollary}

\begin{proof}
Exact observations give $B_0=0$ and the orbit $Z_q$.
Theorem~\ref{thm:v5-dependent} supplies the lower bound
$\epsilon e_M^2(Z_q)$. The resetting suffix observer in
Theorem~\ref{thm:v5-intrinsic} supplies the upper bound $CG_M(q)$
with the same expected losses in the dependent experiment.
The pressure formula and the supremum statement follow from the
corresponding results there. Kernel \eqref{eq:v5-partialreset}
preserves $\mu$ because $T$ does, and has the displayed
minorization. Its transition coin is not an observation.
\end{proof}

\begin{corollary}[Noisy acquisition with retained hidden memory]
\label{cor:v5-dependentnoise}
On $\mathbb T^d$ let $T_bx=bx\pmod1$, with integers $b\ge2,d\ge1$,
let $\mu$ be Haar measure, and use any acquisition kernel $Q$
satisfying \eqref{eq:v5-minorization}. At each acquisition observe
$Y=X+\sigma G\pmod1$, where the observation noises are independent
standard Gaussian vectors. Predict
$\Phi(x)=(\cos(2\pi x_j),\sin(2\pi x_j))_{j=1}^d$.
Assume $W(n)\le Aw_n$ for every $n$ and some $A<\infty$.
With
\[
 a_k=e^{-2\pi^2\sigma^2b^{2k}},\quad v_k=w_ka_k^2,\quad
 B_0=d\sum_{k\ge0}w_k(1-a_k^2),
\]
and $\Psi_v,n_M$ as in \eqref{eq:v4-profile} and
\eqref{eq:v4-state-budget}, for all $M\ge1$ and $\sigma\ge0$,
\begin{equation}\label{eq:v5-dependentnoisybound}
 \epsilon\left(B_0+\frac{\Psi_v(n_M)}{8Ab^6}\right)
 \le R_M^{\mathrm{dep}}
 \le B_0+\frac{d\pi^2}{3}\Psi_v(n_M).
\end{equation}
For geometric acquisitions this also bounds the optimal supremum
raw loss. In particular $R_M^{\mathrm{dep}}\asymp B_0+\Psi_w(n_M)$,
with constants independent of $M$ and $\sigma$ when
$\epsilon$ is bounded below. The critical-window comparison of
Corollary~\ref{cor:v4-critical} remains valid for this dependent
experiment, and the explicit restarting suffix observer has the
same exact crossover coefficient as there.
\end{corollary}

\begin{proof}
The single-acquisition posterior under $\mu K$ is
$m_k(y)=a_k\Phi(b^ky)$, by \eqref{eq:v4-explicit-posterior}.
Lemma~\ref{lem:v4-smallballs} bounds its orbit quantization error
below by $\Psi_v(n_M)/(8Ab^6)$; use
\eqref{eq:v5-dependentlower}. The suffix observer proving
Theorem~\ref{thm:v4-noisy} overwrites its state at acquisitions,
so its exact expected losses, including the formula
\eqref{eq:v4-exact-suffix}, transfer unchanged. This proves the
upper bound and the supremum assertion. Finally
$0\le\Psi_w(n)-\Psi_v(n)\le B_0/d$ proves the raw-profile
comparison. Its constants are uniform in the same compact
$q$-windows as in Corollary~\ref{cor:v4-critical}; the explicit
observer has exactly the same loss, rather than merely comparable
loss, so its coefficient limit transfers as well.
\end{proof}
